
\documentclass{article}
\usepackage{colortbl}
\usepackage{makecell}
\usepackage{multirow}
\usepackage{supertabular}

\begin{document}

\newcounter{utterance}

\centering \large Interaction Transcript for game `function\_detective', experiment `string\_passive\_examples', episode 121 with mistral{-}large{-}3{-}azure.
\vspace{24pt}

{ \footnotesize  \setcounter{utterance}{1}
\setlength{\tabcolsep}{0pt}
\begin{supertabular}{c@{$\;$}|p{.15\linewidth}@{}p{.15\linewidth}p{.15\linewidth}p{.15\linewidth}p{.15\linewidth}p{.15\linewidth}}
    \# & \multicolumn{2}{c}{Player} && \multicolumn{2}{c}{Game Master} \\
    \hline

    \theutterance \stepcounter{utterance}  
    & & & \multicolumn{4}{p{0.6\linewidth}}{
        \cellcolor[rgb]{0.9,0.9,0.9}{
            \makecell[{{p{\linewidth}}}]{
                \texttt{\tiny{[P1$\langle$GM]}}
                \texttt{I have a hidden Python function for you to figure out.} \\
\\ 
\texttt{Your job is to infer what it does and then submit a correct Python implementation.} \\
\\ 
\texttt{Function signature:} \\
\texttt{(text: str) {-}> str} \\
\\ 
\texttt{You have 15 turns total.} \\
\\ 
\texttt{I control which valid input{-}output examples you see. You can request them one at a time.} \\
\\ 
\texttt{You must make your best SOLVE attempt no later than the second{-}to{-}last turn.} \\
\texttt{Do not wait until the very end.} \\
\\ 
\texttt{You can use these actions:} \\
\texttt{{-} JSON format described at the end of the prompt} \\
\texttt{{-} SOLVE: ```python ... ```} \\
\\ 
\texttt{Do not add import statements in your guess.} \\
\texttt{Anything you might need is already available, so just write the function body in hypothesis and return the answer.} \\
\\ 
\texttt{Preloaded examples:} \\
\texttt{1. ('   ', '   ')} \\
\\ 
\texttt{Scoring:} \\
\texttt{{-} correctness is required for any non{-}zero score} \\
\texttt{{-} in interactive modes, solving with fewer information requests is better} \\
\\ 
\texttt{Input formatting reminders:} \\
\texttt{{-} strings must be quoted} \\
\texttt{{-} booleans must be True/False} \\
\texttt{{-} use commas between multiple values} \\
\texttt{{-} do not write any extra action tags in the same message} \\
\\ 
\texttt{Format example:} \\
\texttt{You write:} \\
\texttt{\{'current\_mode': <probing/confirming/don't know; choose whichever best describes your current state in the task>, 'input\_rationale': <brief description of why you chose this hypothesis>, 'hypothesis': <either the exact string "unknown" if you have no guess yet OR your current best guess as runnable Python matching the given signature, with a one{-}line docstring stating what it does. Never write placeholder code, if you cannot commit to a guess, use "unknown". Match the given signature.} \\
\texttt{e.g.} \\
\texttt{def f(x):} \\
\texttt{	"""Returns the absolute value of x."""} \\
\texttt{	return x if x >= 0 else {-}x} \\
\\ 
\texttt{Use "unknown" if you have no guess.>, 'confidence\_score': <score between 1{-}100. 100 means highest confidence in your hypothesis>\}} \\
\\ 
\texttt{Output only one valid json per turn.} \\
            }
        }
    }
     \\ \\

    \theutterance \stepcounter{utterance}  
    & \multicolumn{4}{p{0.6\linewidth}}{
        \cellcolor[rgb]{0.561,0.737,0.561}{
            \makecell[{{p{\linewidth}}}]{
                \texttt{\tiny{[P1$\rangle$GM]}}
                \texttt{```json} \\
\texttt{\{} \\
\texttt{    "current\_mode": "probing",} \\
\texttt{    "input\_rationale": "The preloaded example shows that a string of spaces returns the same string of spaces. This suggests the function might be related to whitespace handling, such as trimming, normalizing, or preserving whitespace. I will request an example with non{-}whitespace characters to narrow down the behavior.",} \\
\texttt{    "hypothesis": "unknown",} \\
\texttt{    "confidence\_score": 10} \\
\texttt{\}} \\
\texttt{```} \\
            }
        }
    }
    & & \\ \\

    \theutterance \stepcounter{utterance}  
    & & & \multicolumn{4}{p{0.6\linewidth}}{
        \cellcolor[rgb]{0.9,0.9,0.9}{
            \makecell[{{p{\linewidth}}}]{
                \texttt{\tiny{[P1$\langle$GM]}}
                \texttt{OUTPUT: ('a{-}b', 'n{-}o')} \\
            }
        }
    }
     \\ \\

    \theutterance \stepcounter{utterance}  
    & \multicolumn{4}{p{0.6\linewidth}}{
        \cellcolor[rgb]{0.561,0.737,0.561}{
            \makecell[{{p{\linewidth}}}]{
                \texttt{\tiny{[P1$\rangle$GM]}}
                \texttt{```json} \\
\texttt{\{} \\
\texttt{    "current\_mode": "probing",} \\
\texttt{    "input\_rationale": "The input 'a{-}b' maps to 'n{-}o'. This suggests a character{-}level transformation, likely a Caesar cipher or similar shift. The letters 'a' and 'n' are 13 positions apart in the alphabet (a=1, n=14), and 'b' and 'o' are also 13 positions apart (b=2, o=15). This is consistent with ROT13, a common cipher that shifts letters by 13 positions. The function likely applies ROT13 to each alphabetic character while leaving non{-}alphabetic characters (like '{-}') unchanged.",} \\
\texttt{    "hypothesis": "def f(text: str) {-}> str:\textbackslash{}n    \textbackslash{}"\textbackslash{}"\textbackslash{}"Applies ROT13 cipher to the input string, shifting alphabetic characters by 13 positions while leaving non{-}alphabetic characters unchanged.\textbackslash{}"\textbackslash{}"\textbackslash{}"\textbackslash{}n    result = {[}{]}\textbackslash{}n    for char in text:\textbackslash{}n        if 'a' <= char <= 'z':\textbackslash{}n            shifted = chr(((ord(char) {-} ord('a') + 13) \% 26) + ord('a'))\textbackslash{}n            result.append(shifted)\textbackslash{}n        elif 'A' <= char <= 'Z':\textbackslash{}n            shifted = chr(((ord(char) {-} ord('A') + 13) \%} \\
            }
        }
    }
    & & \\ \\

    \theutterance \stepcounter{utterance}  
    & & & \multicolumn{4}{p{0.6\linewidth}}{
        \cellcolor[rgb]{0.9,0.9,0.9}{
            \makecell[{{p{\linewidth}}}]{
                \texttt{\tiny{[P1$\langle$GM]}}
                \texttt{Game Violation: ParseError: Response is not valid JSON.} \\
\texttt{Output ONLY:} \\
\texttt{One JSON with confidence\_score, current\_mode, input\_rationale, and hypothesis keys ...} \\
\texttt{OR} \\
\texttt{SOLVE: ```python ... ```} \\
            }
        }
    }
     \\ \\

    \theutterance \stepcounter{utterance}  
    & \multicolumn{4}{p{0.6\linewidth}}{
        \cellcolor[rgb]{0.561,0.737,0.561}{
            \makecell[{{p{\linewidth}}}]{
                \texttt{\tiny{[P1$\rangle$GM]}}
                \texttt{```json} \\
\texttt{\{} \\
\texttt{    "current\_mode": "probing",} \\
\texttt{    "input\_rationale": "The input 'a{-}b' maps to 'n{-}o', which strongly suggests a ROT13 cipher (each letter shifted by 13 positions). The function likely preserves non{-}alphabetic characters like '{-}'. I will now request another example to confirm this hypothesis.",} \\
\texttt{    "hypothesis": "def f(text: str) {-}> str:\textbackslash{}n    \textbackslash{}"\textbackslash{}"\textbackslash{}"Applies ROT13 cipher to the input string, shifting alphabetic characters by 13 positions while leaving non{-}alphabetic characters unchanged.\textbackslash{}"\textbackslash{}"\textbackslash{}"\textbackslash{}n    result = {[}{]}\textbackslash{}n    for char in text:\textbackslash{}n        if 'a' <= char <= 'z':\textbackslash{}n            shifted = chr(((ord(char) {-} ord('a') + 13) \% 26) + ord('a'))\textbackslash{}n            result.append(shifted)\textbackslash{}n        elif 'A' <= char <= 'Z':\textbackslash{}n            shifted = chr(((ord(char) {-} ord('A') + 13) \% 26) + ord('A'))\textbackslash{}n            result.append(shifted)\textbackslash{}n        else:\textbackslash{}n            result.append(char)\textbackslash{}n    return ''.join(result)",} \\
\texttt{    "confidence\_score": 95} \\
\texttt{\}} \\
\texttt{```} \\
            }
        }
    }
    & & \\ \\

    \theutterance \stepcounter{utterance}  
    & & & \multicolumn{4}{p{0.6\linewidth}}{
        \cellcolor[rgb]{0.9,0.9,0.9}{
            \makecell[{{p{\linewidth}}}]{
                \texttt{\tiny{[P1$\langle$GM]}}
                \texttt{OUTPUT: ('a', 'n')} \\
            }
        }
    }
     \\ \\

    \theutterance \stepcounter{utterance}  
    & \multicolumn{4}{p{0.6\linewidth}}{
        \cellcolor[rgb]{0.561,0.737,0.561}{
            \makecell[{{p{\linewidth}}}]{
                \texttt{\tiny{[P1$\rangle$GM]}}
                \texttt{```json} \\
\texttt{\{} \\
\texttt{    "current\_mode": "confirming",} \\
\texttt{    "input\_rationale": "The input 'a' maps to 'n', which is consistent with ROT13 (a→n). This confirms the earlier hypothesis. The function applies ROT13 to alphabetic characters and leaves others unchanged. I will now submit the solution.",} \\
\texttt{    "hypothesis": "def f(text: str) {-}> str:\textbackslash{}n    \textbackslash{}"\textbackslash{}"\textbackslash{}"Applies ROT13 cipher to the input string, shifting alphabetic characters by 13 positions while leaving non{-}alphabetic characters unchanged.\textbackslash{}"\textbackslash{}"\textbackslash{}"\textbackslash{}n    result = {[}{]}\textbackslash{}n    for char in text:\textbackslash{}n        if 'a' <= char <= 'z':\textbackslash{}n            shifted = chr(((ord(char) {-} ord('a') + 13) \% 26) + ord('a'))\textbackslash{}n            result.append(shifted)\textbackslash{}n        elif 'A' <= char <= 'Z':\textbackslash{}n            shifted = chr(((ord(char) {-} ord('A') + 13) \% 26) + ord('A'))\textbackslash{}n            result.append(shifted)\textbackslash{}n        else:\textbackslash{}n            result.append(char)\textbackslash{}n    return ''.join(result)",} \\
\texttt{    "confidence\_score": 100} \\
\texttt{\}} \\
\texttt{```} \\
\\ 
\texttt{SOLVE:} \\
\texttt{```python} \\
\texttt{def f(text: str) {-}> str:} \\
\texttt{    """Applies ROT13 cipher to the} \\
            }
        }
    }
    & & \\ \\

    \theutterance \stepcounter{utterance}  
    & & & \multicolumn{4}{p{0.6\linewidth}}{
        \cellcolor[rgb]{0.9,0.9,0.9}{
            \makecell[{{p{\linewidth}}}]{
                \texttt{\tiny{[P1$\langle$GM]}}
                \texttt{Game Violation: ParseError: Response is not valid JSON.} \\
\texttt{Output ONLY:} \\
\texttt{One JSON with confidence\_score, current\_mode, input\_rationale, and hypothesis keys ...} \\
\texttt{OR} \\
\texttt{SOLVE: ```python ... ```} \\
            }
        }
    }
     \\ \\

    \theutterance \stepcounter{utterance}  
    & \multicolumn{4}{p{0.6\linewidth}}{
        \cellcolor[rgb]{0.561,0.737,0.561}{
            \makecell[{{p{\linewidth}}}]{
                \texttt{\tiny{[P1$\rangle$GM]}}
                \texttt{SOLVE:} \\
\texttt{```python} \\
\texttt{def f(text: str) {-}> str:} \\
\texttt{    """Applies ROT13 cipher to the input string, shifting alphabetic characters by 13 positions while leaving non{-}alphabetic characters unchanged."""} \\
\texttt{    result = {[}{]}} \\
\texttt{    for char in text:} \\
\texttt{        if 'a' <= char <= 'z':} \\
\texttt{            shifted = chr(((ord(char) {-} ord('a') + 13) \% 26) + ord('a'))} \\
\texttt{            result.append(shifted)} \\
\texttt{        elif 'A' <= char <= 'Z':} \\
\texttt{            shifted = chr(((ord(char) {-} ord('A') + 13) \% 26) + ord('A'))} \\
\texttt{            result.append(shifted)} \\
\texttt{        else:} \\
\texttt{            result.append(char)} \\
\texttt{    return ''.join(result)} \\
\texttt{```} \\
            }
        }
    }
    & & \\ \\

    \theutterance \stepcounter{utterance}  
    & & & \multicolumn{2}{p{0.3\linewidth}}{
        \cellcolor[rgb]{0.9,0.9,0.9}{
            \makecell[{{p{\linewidth}}}]{
                \texttt{\tiny{[GM$|$GM]}}
                \texttt{Game Verdict: LOSS} \\
            }
        }
    }
    & & \\ \\

    \theutterance \stepcounter{utterance}  
    & & & \multicolumn{2}{p{0.3\linewidth}}{
        \cellcolor[rgb]{0.9,0.9,0.9}{
            \makecell[{{p{\linewidth}}}]{
                \texttt{\tiny{[GM$|$GM]}}
                \texttt{Incorrect.} \\
\texttt{Failed Cases:} \\
\texttt{Input: {[}'A'{]} | Expected: A | Got: N} \\
\texttt{Input: {[}'aA'{]} | Expected: nA | Got: nN} \\
\texttt{Input: {[}'\_kM'{]} | Expected: \_xM | Got: \_xZ} \\
            }
        }
    }
    & & \\ \\

    \theutterance \stepcounter{utterance}  
    & & & \multicolumn{2}{p{0.3\linewidth}}{
        \cellcolor[rgb]{0.9,0.9,0.9}{
            \makecell[{{p{\linewidth}}}]{
                \texttt{\tiny{[GM$|$GM]}}
                \texttt{def rot13\_lowercase\_only(text: str) {-}> str:} \\
\texttt{    result = ''} \\
\texttt{    for ch in text:} \\
\texttt{        if 'a' <= ch <= 'z':} \\
\texttt{            result += chr((ord(ch) {-} 97 + 13) \% 26 + 97)} \\
\texttt{        else:} \\
\texttt{            result += ch} \\
\texttt{    return result} \\
            }
        }
    }
    & & \\ \\

    \theutterance \stepcounter{utterance}  
    & & & \multicolumn{2}{p{0.3\linewidth}}{
        \cellcolor[rgb]{0.9,0.9,0.9}{
            \makecell[{{p{\linewidth}}}]{
                \texttt{\tiny{[GM$|$GM]}}
                \texttt{def f(text: str) {-}> str:} \\
\texttt{    """Applies ROT13 cipher to the input string, shifting alphabetic characters by 13 positions while leaving non{-}alphabetic characters unchanged."""} \\
\texttt{    result = {[}{]}} \\
\texttt{    for char in text:} \\
\texttt{        if 'a' <= char <= 'z':} \\
\texttt{            shifted = chr(((ord(char) {-} ord('a') + 13) \% 26) + ord('a'))} \\
\texttt{            result.append(shifted)} \\
\texttt{        elif 'A' <= char <= 'Z':} \\
\texttt{            shifted = chr(((ord(char) {-} ord('A') + 13) \% 26) + ord('A'))} \\
\texttt{            result.append(shifted)} \\
\texttt{        else:} \\
\texttt{            result.append(char)} \\
\texttt{    return ''.join(result)} \\
            }
        }
    }
    & & \\ \\

    \theutterance \stepcounter{utterance}  
    & & & \multicolumn{2}{p{0.3\linewidth}}{
        \cellcolor[rgb]{0.9,0.9,0.9}{
            \makecell[{{p{\linewidth}}}]{
                \texttt{\tiny{[GM$|$GM]}}
                \texttt{Mode: passive\_examples} \\
\texttt{Accuracy: 0.230} \\
\texttt{Information requests: 2/15} \\
\texttt{Revealed passive examples: 3} \\
\texttt{Efficiency: 0.933} \\
\texttt{Slack: 13} \\
\texttt{Observed{-}pair consistency: 1.000} \\
\texttt{Parse errors: 2} \\
\texttt{Runtime errors: 0} \\
            }
        }
    }
    & & \\ \\

\end{supertabular}
}

\end{document}
